\section*{Capítulo \ref{ch:g1:landscapes-around-us} --- As paisagens ao nosso redor}

\begin{solution}{exo:g1:landscapes-around-us:1}
\omterm{def:g1:landscapes-around-us:landform}{Montanha}, \omterm{def:g1:landscapes-around-us:landform}{colina}, \omterm{def:g1:landscapes-around-us:landform}{vale}, \omterm{def:g1:landscapes-around-us:landform}{planície}.
\end{solution}

\begin{solution}{exo:g1:landscapes-around-us:2}
Uma \omterm{def:g1:landscapes-around-us:landform}{montanha}: ela é mais alta e mais íngreme que uma \omterm{def:g1:landscapes-around-us:landform}{colina}.
\end{solution}

\begin{solution}{exo:g1:landscapes-around-us:3}
Naturais: o \omterm{def:g1:landscapes-around-us:water}{rio}, a floresta, a praia. Construídos: a estrada, a ponte, a
casa.
\end{solution}

\begin{solution}{exo:g1:landscapes-around-us:4}
Onde a terra encontra o mar.
\end{solution}

\begin{solution}{exo:g1:landscapes-around-us:5}
Por exemplo: natural, as \omterm{def:g1:landscapes-around-us:landform}{montanhas}, a floresta nas encostas ou o riacho;
construído, uma casa ou a igreja da vila.
\end{solution}

\begin{solution}{exo:g1:landscapes-around-us:6}
A \omterm{def:g1:landscapes-around-us:landform}{montanha}, com neve no alto. O \omterm{def:g1:landscapes-around-us:water}{rio} corre no fundo do \omterm{def:g1:landscapes-around-us:landform}{vale}, entre as
\omterm{def:g1:landscapes-around-us:landform}{montanhas} e a \omterm{def:g1:landscapes-around-us:landform}{colina}.
\end{solution}

\begin{solution}{exo:g1:landscapes-around-us:7}
Um \omterm{def:g1:landscapes-around-us:water}{rio} corre de um terreno mais alto (\omterm{def:g1:landscapes-around-us:landform}{colinas} e \omterm{def:g1:landscapes-around-us:landform}{montanhas}) até um \omterm{def:g1:landscapes-around-us:water}{lago} ou
até o mar. A água sempre desce; por isso ela não pode subir até o alto
de uma \omterm{def:g1:landscapes-around-us:landform}{colina}.
\end{solution}

\begin{solution}{exo:g1:landscapes-around-us:8}
A árvore está perto de nós, no primeiro plano; a \omterm{def:g1:landscapes-around-us:landform}{colina} está longe, no
plano de fundo. As coisas que estão longe parecem menores do que as coisas
que estão perto.
\end{solution}

\begin{solution}{exo:g1:landscapes-around-us:9}
Primeiro plano: capim alto e algumas moitas na parte de baixo da foto.
Plano intermediário: o \omterm{def:g1:landscapes-around-us:water}{rio} fazendo curvas pela \omterm{def:g1:landscapes-around-us:landform}{planície}, com árvores ao
longo dele, plantações, um sítio e uma estrada.
Plano de fundo: a terra plana que segue até o horizonte, e o céu.
\end{solution}

\begin{solution}{exo:g1:landscapes-around-us:10}
Porque foram as pessoas que a fizeram: elas limparam o terreno, araram e
semearam o trigo. Sem as pessoas, ali haveria capim, moitas ou floresta, e
não uma plantação de trigo.
\end{solution}

\begin{solution}{exo:g1:landscapes-around-us:11}
Não mudaram: a \omterm{def:g1:landscapes-around-us:landform}{colina} e o \omterm{def:g1:landscapes-around-us:water}{rio}. Acrescentados pelas pessoas: as
plantações, mais casas, a estrada e a ponte (e a maior parte da floresta
foi derrubada).
\end{solution}
